\section*{अध्याय \ref{ch:b3:bacteriology} --- जीवाणु-विज्ञान: वृद्धि, शरीरक्रिया और आनुवंशिकी}

\begin{solution}{exo:b3:bacteriology:1}
\omterm{def:b3:bacteriology:envelope}{ग्राम-धनात्मक}: कोशिकाद्रव्यी झिल्ली के चारों ओर मोटी, बहुपरती
\omterm{def:b3:bacteriology:envelope}{पेप्टिडोग्लाइकन} दीवार जिसमें टाइकोइक अम्ल पिरोए हैं। \omterm{def:b3:bacteriology:envelope}{ग्राम-ऋणात्मक}:
किसी परिकोशिकाद्रव्य में पतली \omterm{def:b3:bacteriology:envelope}{पेप्टिडोग्लाइकन} परत, फिर कोई \omterm{def:b3:bacteriology:envelope}{बाह्य
झिल्ली} जिसका बाहरी पटल \omterm{def:b3:bacteriology:envelope}{लिपोपॉलिसैकेराइड} है, और पोरिन। क्रिस्टल
वायलेट--आयोडीन संकुल \omterm{def:b3:bacteriology:envelope}{ग्राम-धनात्मक} की मोटी, ऐल्कोहल-निर्जलित दीवार में
फँसा रह जाता है (बैंगनी) पर \omterm{def:b3:bacteriology:envelope}{ग्राम-ऋणात्मक} की पतली दीवार और अस्तव्यस्त
\omterm{def:b3:bacteriology:envelope}{बाह्य झिल्ली} से धुल जाता है, और वह तब गुलाबी प्रति-अभिरंजन ले लेती है।
\end{solution}

\begin{solution}{exo:b3:bacteriology:2}
$10^{6}$ गुना $= 2^{19.9}$: लगभग $20$ पीढ़ियाँ; $T = \qty{8}{h}/20
= \qty{24}{min}$; $\mu = \ln 2/\qty{0.4}{h} = \qty{1.7}{h^{-1}}$।
\end{solution}

\begin{solution}{exo:b3:bacteriology:3}
\omterm{def:b3:bacteriology:hgt}{रूपांतरण} --- ग्रिफ़िथ (1928) और एवरी, मैक्लॉड तथा मैक्कार्टी
(1944): ताप से मारे गए विषालु न्यूमोकोकस जीवित अविषालु न्यूमोकोकस को
रूपांतरित कर देते हैं, और वह कारक DNA है। \omterm{def:b3:bacteriology:hgt}{संचारण} --- ज़िंडर और लेडरबर्ग
(1952): फ़ेज P22 किसी ऐसे छन्ने के पार, जो कोशिकाएँ रोक देता है, एक
प्रजाति से दूसरी तक \emph{Salmonella} जीन ले जाता है। \omterm{def:b3:bacteriology:hgt}{संयुग्मन} ---
लेडरबर्ग और टेटम (1946): \emph{E.~coli} की दो पोषकमाँगी प्रजातियाँ
मिलाने पर ही पूर्णपोषी पुनर्योगज देती हैं, और इसके लिए कोशिका-संपर्क
चाहिए (डेविस की U-नली)।
\end{solution}

\begin{solution}{exo:b3:bacteriology:4}
पेनिसिलिन: वे ट्रांसपेप्टिडेज़ जो \omterm{def:b3:bacteriology:envelope}{पेप्टिडोग्लाइकन} को तिर्यक-बद्ध करती
हैं, जो पशु-कोशिकाओं में हैं ही नहीं। स्ट्रेप्टोमाइसिन: लघु राइबोसोमी
उपइकाई, जिसका जीवाणु रूप सुकेंद्रकी रूप से भिन्न है। सिप्रोफ़्लोक्सैसिन:
DNA गाइरेज़, यानी कोई जीवाणु सांस्थितिकी-एंज़ाइम जो मनुष्यों में नहीं
है। रिफ़ैम्पिसिन: जीवाणु RNA पॉलीमरेज़ की $\beta$ उपइकाई, जो सुकेंद्रकी
एंज़ाइम में नहीं है।
\end{solution}

\begin{solution}{exo:b3:bacteriology:5}
$S^{*} = K_{S}D/(\mu_{\max} - D)$, $X^{*} = Y(S_{0} - S^{*})$। $D = 0.2$:
$S^{*} = \qty{0.0067}{g/L}$, $X^{*} = \qty{0.80}{g/L}$। $D = 0.6$:
$S^{*} = 0.06$, $X^{*} = 0.78$। $D = 0.79$: $S^{*} = 1.58$, $X^{*} = 0.17$।
बहना तब जब $D = \mu(S_{0}) = 0.8\times 2/2.02 = \qty{0.79}{h^{-1}}$ हो।
\end{solution}

\begin{solution}{exo:b3:bacteriology:6}
$(N/V)_{\text{देहली}} = kA_{\text{देहली}}/p$। प्रकाश-अंग: $0.01\times
6\times 10^{12}/100 = 6\times 10^{8}$ प्रति लीटर $= 6\times 10^{5}$ प्रति
mL। समुद्री जल: $10\times 6\times 10^{12}/100 = 6\times 10^{11}$ प्रति लीटर
$= 6\times 10^{8}$ प्रति mL --- हज़ार गुना ऊँचा, जो समुद्र में कभी
नहीं आता।
\end{solution}

\begin{solution}{exo:b3:bacteriology:7}
\qty{25}{mm}: संवेदी, नीचा MIC। \qty{12}{mm}: घटी हुई संवेदिता,
मध्यवर्ती --- साधारण मात्राओं पर विफल हो सकती है। \qty{0}{mm}: प्रतिरोधी,
चक्रिका तक वृद्धि। औषध चक्रिका से इस तरह विसरित होती है कि उसकी
सांद्रता दूरी के साथ किसी पुनरुत्पाद्य ढंग से गिरती है; वृद्धि वहाँ रुक
जाती है जहाँ सांद्रता MIC के बराबर हो, इसलिए क्षेत्र-त्रिज्या और MIC
ज्ञात MIC वाली प्रजातियों से बने किसी अंशांकन वक्र से जुड़े हैं।
\end{solution}

\begin{solution}{exo:b3:bacteriology:8}
संयुग्मी \omterm{def:b3:genetic-engineering:tools}{प्लाज़्मिड} कोई ऐसा इंटीग्रॉन ढोता है जिसकी पेटी-शृंखला में कई
प्रतिरोध-जीन हैं, और साथ में और ढोते ट्रांसपोज़ॉन; \omterm{def:b3:genetic-engineering:tools}{प्लाज़्मिड} एक ही
संगम में पार चला जाता है और उन सबको अभिव्यक्त करता है। वे इसलिए गुच्छित
रहते हैं कि इनमें से किसी भी एक औषध का चयन पूरा \omterm{def:b3:genetic-engineering:tools}{प्लाज़्मिड} बनाए रखता
है, इसलिए उस पर सवारी करते जीन टिके रहते हैं (सह-चयन); इंटीग्रॉन
पेटियाँ पकड़ लेते हैं; ट्रांसपोज़ॉन \omterm{def:b3:genetic-engineering:tools}{प्लाज़्मिडों} पर कूद जाते हैं; और
स्थानांतरण की इकाई जीन नहीं, \omterm{def:b3:genetic-engineering:tools}{प्लाज़्मिड} है।
\end{solution}

\begin{solution}{exo:b3:bacteriology:9}
अकेली औषध: प्रत्याशित $10^{-7}\times 10^{9} = 100$ प्रतिरोधी कोशिकाएँ। दो
स्वतंत्र औषधें: $10^{-14}\times 10^{9} = 10^{-5}$। अटल कोशिकाएँ
उत्परिवर्ती नहीं हैं: सुप्त कोशिकाओं का कोई छोटा अंश किसी भी औषध से बच
जाता है और औषध हटते ही किसी संवेदी समष्टि में फिर उग आता है, इसलिए कोई
दूसरी औषध मदद नहीं करती; केवल इतना लंबा उपचार, जो उन्हें जागते समय
पकड़ ले, या सुप्त कोशिकाओं पर सक्रिय कोई औषध ही करती है।
\end{solution}

\begin{solution}{exo:b3:bacteriology:10}
$X > X^{*}$: खपत $\mu X/Y$ आपूर्ति से अधिक है, इसलिए $S$ $S^{*}$ से
नीचे गिरता है; तब $\mu(S) < D$ और $\mathrm{d}X/\mathrm{d}t < 0$, $X$
वापस गिरता है; जैसे-जैसे $X$ गिरता है, $S$ फिर चढ़ता है। दोनों विचलन
सुधर जाते हैं: स्थायी अवस्था स्थिर है। दो जातियाँ: टिके रहने के लिए हर
एक को $\mu_{i}(S) = D$ चाहिए, जो उसका अपना $S_{i}^{*}$ तय कर देता है;
पोषक केवल एक मान ले सकता है, इसलिए अधिक-से-अधिक एक जाति टिकती है। जिसका
$S^{*}(D)$ नीचा हो वह $S$ को उससे नीचे गिरा देती है जो दूसरी को चाहिए,
और दूसरी बह जाती है; वह जाति कौन-सी होगी यह $D$ के साथ बदल सकता है यदि
उनके मोनो वक्र कटते हों।
\end{solution}

\begin{solution}{exo:b3:bacteriology:11}
घनत्व $\rho = N/V$ के साथ, बंद अवस्था $A = p_{0}\rho/k$ तब तक रहती है
जब तक वह $A_{\text{देहली}}$ से नीचे रहे, यानी $\rho < kA_{\text{देहली}}/p_{0}$;
चालू अवस्था $A = p_{1}\rho/k$ तब तक रहती है जब तक वह ऊपर रहे, यानी $\rho
> kA_{\text{देहली}}/p_{1}$। $kA_{\text{देहली}}/p_{1} < \rho <
kA_{\text{देहली}}/p_{0}$ के लिए दोनों स्व-संगत हैं: जो समष्टि चालू हो
चुकी है वह अपने ऊँचे उत्पादन से स्वयं को चालू रखती है, और जो नहीं हुई
वह बंद रहती है। उपयोग: प्रतिबद्धता --- एक बार कोई बस्ती कोई जैवफ़िल्म
बनाने या विष छोड़ने का निर्णय कर ले, तो घनत्व में कोई गिरावट वह निर्णय
नहीं पलटती, और स्विच श्रेणीबद्ध के बजाय तीखा तथा एक साथ होता है।
\end{solution}

\begin{solution}{exo:b3:bacteriology:12}
(क) छोटे कोर्स में ऊँची मात्राएँ: स्तर उत्परिवर्ती MIC से ऊपर बना रहता
है और वन्य-प्ररूप तथा प्रथम-चरण उत्परिवर्ती दोनों को मारता है ---
पक्ष में, जिसकी सीमा विषालुता है। (ख) लंबे कोर्स में कम मात्राएँ: स्तर
दिनों तक खिड़की में बैठा रहता है और प्रतिरोध पनपाता है --- विपक्ष में।
(ग) लक्षण मिटते ही रोक देना: जो बचते हैं वे सबसे कम संवेदी कोशिकाएँ हैं
और गिरता स्तर खिड़की से होकर गुज़रता है --- विपक्ष में। (घ) संयोजन:
किसी कोशिका को दोनों के प्रति प्रतिरोधी होना पड़ेगा, जो एक अरब कोशिकाओं
में नहीं होता --- पक्ष में।
\end{solution}

\begin{solution}{pb:b3:bacteriology:1}
\textbf{1.} $18$ पीढ़ियाँ: $2^{18} = 2.6\times 10^{5}$ कोशिकाएँ,
प्रति mL $2.6\times 10^{3}$।
\textbf{2.} $2\times 10^{9}\times 100 = 2\times 10^{11}$ कोशिकाएँ $=
2^{37.5}$: \qty{12.5}{h} बाद; द्रव्यमान \qty{0.2}{g}।
\textbf{3.} $6\times 10^{27}/10^{-12} = 6\times 10^{39} = 2^{132}$ कोशिकाएँ:
$132$ पीढ़ियाँ, \qty{44}{h}। पोषक, ऑक्सीजन और अपशिष्ट उसे एक दिन के
भीतर रोक देते हैं; चरघातांकी वृद्धि सदा अल्पकालिक होती है।
\textbf{4.} स्थिर कोशिकाएँ अपने राइबोसोम उतार चुकी होती हैं और अपने
उपापचयी एंज़ाइम दमित कर चुकी होती हैं; विभाजित होने से पहले उन्हें वे
फिर बनाने पड़ते हैं। चरघातांकी कोशिकाएँ पूरी तरह सज्जित आती हैं।
\textbf{5.} $\mu = \ln 2/\qty{0.333}{h} = \qty{2.08}{h^{-1}}$; $T =
\qty{60}{min}$ पर \qty{0.69}{h^{-1}}; \qty{6}{h} बाद $2^{6} = 64$,
$2^{18}$ के बजाय: गुणक $2^{12} \approx 4000$ कम।
\textbf{6.} प्रति mL $200/(0.1\times 10^{-6}) = 2\times 10^{9}$ जीवक्षम
कोशिकाएँ। सूक्ष्मदर्शी मरी कोशिकाएँ और वे भी गिन लेता है जो पट्टिका पर
नहीं उगेंगी, और कोई गुच्छा एक ही बस्ती देता है।
\textbf{7.} $D = 0.5/1 = \qty{0.5}{h^{-1}}$; $T = \ln 2/0.5 =
\qty{1.39}{h} = \qty{83}{min}$।
\textbf{8.} $S^{*} = 0.01\times 0.5/0.5 = \qty{0.01}{g/L}$; $X^{*} =
0.5\times 1.99 \approx \qty{1}{g/L} = 10^{12}$ कोशिकाएँ प्रति लीटर,
प्रति mL $10^{9}$।
\textbf{9.} प्रति घंटा $0.5\times 10^{12} = 5\times 10^{11}$ कोशिकाएँ;
महीने में $720/1.39
\approx 520$ पीढ़ियाँ।
\textbf{10.} $D = 0.95$: $S^{*} = 0.01\times 0.95/0.05 = \qty{0.19}{g/L}$,
$X^{*} = 0.5\times 1.81 = \qty{0.9}{g/L}$। $F = \qty{1.1}{L/h}$ पर $D
> \mu_{\max}$: बह जाना।
\textbf{11.} उत्परिवर्ती $S^{*} = 0.01\times 0.5/0.7 = \qty{0.0071}{g/L}$
पर रखता है; उस ग्लूकोज़ पर मूल प्रजाति $1.0\times 0.0071/0.0171 =
\qty{0.42}{h^{-1}} < D$ से बढ़ती है और $e^{-0.08t}$ की तरह घटती है: कुछ
हफ़्तों में बह जाती है। उत्परिवर्ती छा जाता है।
\textbf{12.} प्रति पीढ़ी $10^{-9}\times 10^{12} = 1000$ लाभकारी
उत्परिवर्तन: अनुकूलन सतत है, और हर कुछ सौ पीढ़ियों में कोई नया लाभकारी
क्लोन छा जाता है --- केमोस्टैट एक विकास-मशीन है।
\textbf{13.} $kA_{\text{देहली}}/p = 0.05\times 6\times 10^{12}/500 =
6\times 10^{8}$ प्रति लीटर $= 6\times 10^{5}$ कोशिकाएँ प्रति mL।
\textbf{14.} $10^{10}/6\times 10^{5} \approx 1.7\times 10^{4}$ गुना।
\textbf{15.} $20\times 6\times 10^{12}/500 = 2.4\times 10^{11}$ प्रति लीटर,
प्रति mL $2.4\times 10^{8}$ --- समुद्र की प्रति mL $10^{2}$ से छह कोटि
ऊपर: अँधेरा।
\textbf{16.} दस गुना यानी $3.3$ द्विगुणन, \qty{6.6}{h}। प्रति mL $10^{9}$
पर समष्टि अब भी देहली से $1700$ गुना ऊपर है: प्रकाश जलता रहता है।
\textbf{17.} समुद्र में प्रकाश कोई नहीं देखता और वह किसी परपोषी को नहीं
खिलाता: कोई अँधेरा उत्परिवर्ती \qty{20}{\%} बचाकर चमकनेवालों से आगे
बढ़ जाता है। अंग में स्क्विड प्रकाश का चयन करता है --- वह अँधेरे
धोखेबाज़ों को निकाल देता है या उन्हें नहीं खिलाता --- इसलिए वह लागत एक
घर खरीद देती है।
\textbf{18.} ग्राही रोकिए या \omterm{def:b3:bacteriology:quorum}{स्वप्रेरक} नष्ट कीजिए (गणपूर्ति-शमन):
जीवाणु जीते रहते हैं पर अपने विष कभी नहीं छोड़ते, और प्रतिरक्षा-तंत्र
उन्हें साफ़ कर देता है। चूँकि संवेदी और ``प्रतिरोधी'' जीवाणु बराबर बढ़ते
हैं --- कुछ मारा ही नहीं जाता --- औषध से बच निकलने का वरणीय लाभ छोटा
है, और प्रतिरोध किसी मारने वाली औषध की तुलना में धीरे फैलना चाहिए।
\textbf{19.} A: $10^{-8}\times 10^{9} = 10$; B: $100$; दोनों: $10^{-15}\times
10^{9} = 10^{-6}$।
\textbf{20.} $P_{0}(\text{A}) = e^{-10} = 5\times 10^{-5}$;
$P_{0}(\text{दोनों}) = e^{-10^{-6}} = 0.999999$।
\textbf{21.} $16 \to 8$: एक अर्ध-आयु, \qty{4}{h}; $16 \to 1$: चार,
\qty{16}{h}; \qty{4}{h} से \qty{16}{h} तक खिड़की में, यानी प्रति मात्रा
\qty{12}{h}।
\textbf{22.} हर \qty{12}{h} में मात्रा: घंटे $4$ से घंटे
$12$ तक खिड़की में, यानी दो-तिहाई समय; दस दिनों में A के प्रथम-चरण
उत्परिवर्ती --- आरंभ में उनमें से दस --- चुन लिए जाते हैं और अकेली A
विफल हो जाती है।
\textbf{23.} पूरे समय \qty{8}{mg/L} से ऊपर: हर अर्ध-आयु पर मात्रा
(\qty{4}{h}), या हर \qty{12}{h} में \qty{64}{mg/L} पर शिखर छूती कोई मात्रा
($64 \to 8$ तीन अर्ध-आयुओं में), या कोई सतत आसेचन। लागतें: ऊँचे शिखरों
की विषालुता, या दिन में छह मात्राएँ जो रोगी निभाएँगे नहीं।
\textbf{24.} $10^{-4}\times 10^{9} = 10^{5}$ अटल कोशिकाएँ औषधें जो भी
हों बच जाती हैं; उपचार रुकते ही वे जागकर कोई पूरी तरह संवेदी समष्टि
फिर उगा देती हैं --- कोई पुनरावर्तन। उपचार इतना लंबा चलना चाहिए कि
उन्हें सुप्तावस्था छोड़ते समय पकड़ा जा सके, इसीलिए यक्ष्मा छह महीने लेती है।
\textbf{25.} केमोस्टैट में प्रति mL $10^{9}$ कोशिकाएँ; प्रकाश प्रति mL
$6\times 10^{5}$ कोशिकाओं पर चालू होता है; $P_{0} = 0.999999$ कि
दो-औषध कोर्स को कोई प्रतिरोधी कोशिका न मिले।
\end{solution}
